\chapterbackmatter{Solutions to Supplementary Exercises}

\begin{enumerate}
  \SupplementarySolution{1}{The Free-Scalar Short-Distance Product}
  Wick's theorem gives
  \[
    \phi(x)\phi(0)=D(x)\mathbf1+:\!\phi(x)\phi(0)\!:,
    \qquad D(x)=\frac1{4\pi^2x^2}.
  \]
  Taylor expansion inside normal ordering yields
  \[
    :\!\phi(x)\phi(0)\!:
      =:\!\phi^2(0)\!:
       +x^\mu:\!\partial_\mu\phi(0)\phi(0)\!:+O(x^2).
  \]
  Since
  \(2:\!\phi\partial_\mu\phi\!:
    =\partial_\mu:\!\phi^2\!:\), the result is
  \[
    \boxed{\phi(x)\phi(0)
      =\frac{\mathbf1}{4\pi^2x^2}
       +:\!\phi^2(0)\!:
       +\frac{x^\mu}{2}\partial_\mu:\!\phi^2(0)\!:
       +O(x^2).}
  \]
  The same coefficient multiplies every matrix element, illustrating the
  state independence of a Wilson coefficient.

  \SupplementarySolution{2}{A Normal-Ordered Square OPE}
  Only cross-contractions are allowed between the two normal-ordered
  factors.  There are two double contractions and four single
  contractions:
  \[
    O(x)O(0)
      =2D(x)^2\mathbf1
       +4D(x):\!\phi(x)\phi(0)\!:
       +:\!\phi^2(x)\phi^2(0)\!:.
  \]
  The last term begins with the dimension-four operator
  \(:\!\phi^4\!:\).  Keeping the requested terms and using the expansion
  from Solution 1 gives
  \[
    \boxed{O(x)O(0)
      =2D^2\mathbf1+4D:\!\phi^2\!:
       +2D\,x^\mu\partial_\mu:\!\phi^2\!:+O(d\geq4).}
  \]
  Thus the descendant coefficient is fixed by translation invariance
  once the primary coefficient is known.

  \SupplementarySolution{3}{Dimension-Four Operators in a Free OPE}
  If \(k\) fields are contracted across the two factors, the
  combinatorial coefficient is
  \(\binom4k^2k!\), and the uncontracted operator is
  \(:\!\phi^{8-2k}\!:\).  Hence
  \[
    :\!\phi^4(x)\!::\!\phi^4(0)\!:
      =\sum_{k=0}^4\binom4k^2k!D(x)^k
        :\!\phi^{4-k}(x)\phi^{4-k}(0)\!:.
  \]
  The operators of dimension at most four arise for \(k=4,3,2\):
  \[
    \boxed{24D^4\mathbf1
      +96D^3:\!\phi^2\!:
      +72D^2:\!\phi^4\!:+\text{descendants and }d>4.}
  \]
  Taylor expansion of the uncontracted fields supplies descendants; it
  does not change these leading primary coefficients.

  \SupplementarySolution{4}{A Wilson Coefficient from Soft-State Matching}
  The vacuum matrix element gives
  \[
    \langle\Omega|A(x)B(0)|\Omega\rangle
      =C_1(x)+C_{\phi^2}(x)
       \langle\Omega|[\phi^2]_R|\Omega\rangle.
  \]
  In a normal-ordering or zero-condensate matching scheme this directly
  fixes \(C_1\).  Next take an amputated two-particle matrix element:
  \[
    \mathcal G_2(x;p_1,p_2)
      =C_1\langle p_1p_2|\mathbf1|0\rangle
       +C_{\phi^2}
        \langle p_1p_2|[\phi^2]_R|0\rangle.
  \]
  One must subtract the already determined identity contribution,
  including disconnected pieces and operator-mixing counterterms.
  Dividing the remainder by the renormalized \(\phi^2\) matrix element
  determines \(C_{\phi^2}\).  Generic off-shell soft momenta separate
  ultraviolet matching from infrared effects.

  \SupplementarySolution{5}{Beta Functions under a Coupling Redefinition}
  The chain rule gives the exact transformation law
  \[
    \boxed{\bar\beta(\bar g)
      =\frac{d\bar g}{dg}\,\beta(g).}
  \]
  Inverting the redefinition,
  \[
    g=\bar g-a_2\bar g^2+(2a_2^2-a_3)\bar g^3+\cdots,
  \]
  and re-expanding yields
  \[
    \bar\beta(\bar g)
      =b_2\bar g^2+b_3\bar g^3
       +\bigl[b_4-a_2b_3+(a_3-a_2^2)b_2\bigr]\bar g^4+\cdots.
  \]
  Thus \(b_2\) and \(b_3\) are invariant, whereas \(b_4\) is
  scheme-dependent.  Geometrically, the beta function is a vector field
  on the one-dimensional space of couplings.

  \SupplementarySolution{6}{Dimensional Transmutation}
  Separating variables gives
  \[
    \frac1{g(\mu)}=b\ln\frac{\mu}{\Lambda},\qquad
    \boxed{g(\mu)=\frac1{b\ln(\mu/\Lambda)}}.
  \]
  Equivalently,
  \[
    \Lambda=\mu\exp\!\left[-\frac1{b\,g(\mu)}\right].
  \]
  Differentiation using the beta function shows
  \(\mu\,d\Lambda/d\mu=0\).  Specifying \(g\) at one scale is therefore
  equivalent to specifying the dimensionful constant \(\Lambda\).
  Classical scale invariance is replaced by dependence on this
  dynamically generated scale; this is dimensional transmutation.

  \SupplementarySolution{7}{Transversality of a Current Correlator}
  Define
  \[
    \Pi^{\mu\nu}(q)
      =i\int d^4x\,e^{-iq\cdot x}
       \langle\Omega|T\{J^\mu(x)J^\nu(0)\}|\Omega\rangle.
  \]
  Current conservation and the Ward identity give
  \(q_\mu\Pi^{\mu\nu}=0\), apart from local contact polynomials.
  Lorentz covariance then leaves one nonlocal tensor:
  \[
    \boxed{\Pi^{\mu\nu}
      =(q^\mu q^\nu-q^2\eta^{\mu\nu})\Pi(q^2)
       +\text{contact terms}.}
  \]
  Inserting a complete set of states shows that the transverse spectral
  density is nonnegative for a Hermitian current: every contraction with
  a physical transverse polarization is a sum of absolute squares.

  \SupplementarySolution{8}{Large-Momentum Expansion of a Dispersion Relation}
  For \(Q^2\) larger than the relevant spectral scales,
  \[
    \frac1{s+Q^2}
      =\frac1{Q^2}\left(1-\frac{s}{Q^2}
        +\frac{s^2}{Q^4}-\cdots\right).
  \]
  Therefore
  \[
    \boxed{\Pi(Q^2)
      =\frac{M_0}{Q^2}-\frac{M_1}{Q^4}
       +\frac{M_2}{Q^6}+\cdots,\qquad
       M_n=\int_0^\infty ds\,s^n\rho(s).}
  \]
  The interchange of series and integral requires the displayed moments
  to converge.  In a field theory they often do not; subtractions,
  weighted differences of correlators, or an asymptotic OPE must then be
  used instead of treating the expansion as a convergent identity.

  \SupplementarySolution{9}{A Finite-Energy Sum Rule}
  Let
  \(\rho(s)=\operatorname{Im}\Pi(s+i0)/\pi\).
  Deforming the counterclockwise contour around the cut gives
  \[
    \boxed{\int_0^{s_0}ds\,s^n\rho(s)
      =-\frac1{2\pi i}\oint_{|s|=s_0}ds\,s^n\Pi(s).}
  \]
  The sign follows because the two sides of the cut are traversed
  oppositely when the outer circle is counterclockwise.  If \(s_0\) is
  large, the correlator on the circle can be replaced by its OPE.  The
  equality then relates a measured spectral moment to perturbative
  terms and vacuum condensates.

  \SupplementarySolution{10}{Borel Suppression of a Continuum}
  With the standard QCD-sum-rule definition
  \[
    \mathcal B_{M^2}
      =\lim_{\substack{Q^2,n\to\infty\\Q^2/n=M^2}}
       \frac{(Q^2)^n}{(n-1)!}
       \left(-\frac{d}{dQ^2}\right)^n,
  \]
  one finds
  \[
    \boxed{\mathcal B_{M^2}\frac1{s+Q^2}
      =\frac1{M^2}e^{-s/M^2}.}
  \]
  Thus the dispersion relation becomes exponentially weighted, so
  states with \(s\gg M^2\) are suppressed relative to a low resonance.
  Any finite polynomial in \(Q^2\) is annihilated after sufficiently many
  derivatives, which removes the subtraction ambiguity.

  \SupplementarySolution{11}{A Pole-plus-Continuum Spectral Model}
  The Borel-transformed dispersion relation is
  \[
    \mathcal B\Pi
      =\frac1{M^2}\left[
        f_R^2e^{-m_R^2/M^2}
        +\int_{s_0}^\infty ds\,
          \rho_{\rm pert}(s)e^{-s/M^2}\right].
  \]
  After moving the modeled continuum to the OPE side, define
  \[
    R_0(M^2)=f_R^2e^{-m_R^2/M^2},\qquad
    R_1(M^2)=-\frac{dR_0}{d(1/M^2)}
      =m_R^2R_0.
  \]
  Hence
  \[
    \boxed{m_R^2=\frac{R_1(M^2)}{R_0(M^2)}.}
  \]
  A reliable extraction requires a window in which the truncated OPE
  converges and the pole is not overwhelmed by the continuum.

  \SupplementarySolution{12}{Condensate Dimension and Power Suppression}
  Since the invariant amplitude is dimensionless, an operator
  \(O_d\) of dimension \(d\) must occur as
  \[
    \Pi(Q^2)\supset
      \frac{C_d(Q^2/\mu^2,\alpha_s)}{Q^d}
      \langle O_d(\mu)\rangle,
  \]
  apart from logarithms and possible powers already removed in the
  tensor decomposition.  Thus the gluon condensate contributes as
  \(\langle F^2\rangle/Q^4\), whereas a four-quark condensate contributes
  as \(\langle(\bar q\Gamma q)^2\rangle/Q^6\).  The latter is suppressed
  by an additional \(1/Q^2\), although its numerical matrix element and
  Wilson coefficient also matter.

  \SupplementarySolution{13}{The Inclusive Hadronic Tensor}
  A symmetric tensor made from \(p,q,\eta\) initially has four
  independent structures after identifying the equal mixed
  coefficients.  The two equations from
  \(q_\mu W^{\mu\nu}=0\) leave two independent combinations.  A convenient
  basis is
  \[
    \boxed{\begin{aligned}
    W^{\mu\nu}={}&-\left(\frac{q^\mu q^\nu}{q^2}
      -\eta^{\mu\nu}\right)W_1\\
      &+\frac1{m_N^2}
      \left(p^\mu-\frac{p\cdot q}{q^2}q^\mu\right)
      \left(p^\nu-\frac{p\cdot q}{q^2}q^\nu\right)W_2 .
    \end{aligned}}
  \]
  Both bracketed tensors are transverse.  Parity violation would allow
  one additional antisymmetric structure proportional to
  \(\epsilon^{\mu\nu\rho\sigma}p_\rho q_\sigma\).

  \SupplementarySolution{14}{Bjorken Support from Parton Kinematics}
  With the project's mostly-plus metric, put \(Q^2=q^2>0\) for a
  spacelike transfer.  Neglecting the target and parton masses,
  \[
    0=(xp+q)^2=Q^2+2x\,p\cdot q.
  \]
  Therefore
  \[
    \boxed{x=x_B=-\frac{Q^2}{2p\cdot q}.}
  \]
  A parton moving with the parent hadron carries a positive fraction of
  its longitudinal momentum, so \(x>0\); it cannot carry more momentum
  than the whole hadron, so \(x\leq1\).  The endpoint \(x=1\) is the
  elastic parton limit, and \(x_B>1\) has no massless-parton support.

  \SupplementarySolution{15}{Parton-Model Structure Functions}
  The incoherent probability to strike a quark or antiquark of flavor
  \(f\) is weighted by the squared electromagnetic charge.  The
  on-shell delta function fixes the parton fraction to Bjorken \(x\),
  giving
  \[
    \boxed{F_2(x)=x\sum_fQ_f^2
      \bigl[q_f(x)+\bar q_f(x)\bigr].}
  \]
  The transverse response of a massless Dirac constituent gives the
  Callan--Gross relation,
  \[
    \boxed{F_1(x)=\frac{F_2(x)}{2x}
      =\frac12\sum_fQ_f^2
       [q_f(x)+\bar q_f(x)].}
  \]
  Gluons enter these electromagnetic functions through QCD corrections,
  not as directly charged leading-order targets.

  \SupplementarySolution{16}{The Callan--Gross Diagnostic}
  The pointlike Dirac tensor has only transverse absorption at leading
  order.  Using \(2xF_1=F_2\),
  \[
    \boxed{F_L(x)\equiv F_2(x)-2xF_1(x)=0.}
  \]
  A scalar charged parton would have a different longitudinal response,
  so a leading, unsuppressed violation would test the constituent spin.
  In QCD, however, gluon radiation and quark transverse momentum generate
  \(F_L\neq0\) at higher orders even though the quarks have spin
  \(\tfrac12\).  Target-mass and higher-twist corrections also violate
  the ideal relation by powers of \(1/Q\).

  \SupplementarySolution{17}{Even Spins in the Forward OPE}
  Crossing the two electromagnetic currents exchanges absorption and
  emission and sends \(\nu\to-\nu\), hence
  \(\omega\to-\omega\).  The spin-averaged parity-even forward amplitude
  is symmetric:
  \[
    T_r(\omega,Q^2)=T_r(-\omega,Q^2).
  \]
  Its analytic expansion below threshold is therefore
  \[
    T_r=\sum_{n=0}^\infty c_{2n}(Q^2)\omega^{2n}.
  \]
  A rank-\(s\) symmetric traceless operator contributes \(s\) powers of
  the target momentum and hence a term proportional to \(\omega^s\).
  Only even-spin operators occur in this channel.  Odd spins appear in
  crossing-odd or parity-violating combinations.

  \SupplementarySolution{18}{Moments and Twist}
  The matrix element of a spin-\(s\) operator supplies \(s\) powers of the
  target momentum.  Its Wilson coefficient must compensate the
  dimension \(d\).  At fixed Bjorken variable, the net hard-scale factor
  is
  \[
    (Q^2)^{(2-\tau)/2}=Q^{2-\tau},
    \qquad \tau=d-s.
  \]
  Relative to the leading twist-two contribution,
  \[
    \boxed{\frac{T_\tau}{T_{\tau=2}}
      \sim\frac1{Q^{\tau-2}}}
  \]
  up to anomalous-dimension logarithms.  Thus twist four gives a
  \(1/Q^2\) correction, while twist six gives \(1/Q^4\).

  \SupplementarySolution{19}{Valence-Number Sum Rules}
  The proton contains two net up quarks, one net down quark, and no net
  strange or heavier flavor.  Therefore
  \[
    \boxed{\int_0^1u_v(x)\,dx=2,\qquad
      \int_0^1d_v(x)\,dx=1,\qquad
      \int_0^1q_v(x)\,dx=0\quad(q=s,c,b,\ldots).}
  \]
  Sea pairs contribute equally to \(q\) and \(\bar q\) and cancel in
  \(q_v\).  In the OPE these zeroth moments are forward matrix elements
  of conserved flavor-vector currents.  Their zero anomalous dimensions
  protect the sum rules under QCD evolution.

  \SupplementarySolution{20}{The Momentum Sum Rule}
  At every factorization scale,
  \[
    \boxed{\int_0^1dx\,x\left[
      g(x,\mu)+\sum_f(q_f(x,\mu)+\bar q_f(x,\mu))\right]=1.}
  \]
  Splittings continually transfer momentum between quarks and gluons, so
  none of the separate integrals is scale independent.  Their sum is the
  forward matrix element of the conserved energy-momentum tensor.  The
  corresponding spin-two singlet anomalous-dimension matrix has a zero
  eigenvalue, ensuring that DGLAP evolution preserves the total.

  \SupplementarySolution{21}{Testing a Plus Distribution}
  For the zeroth moment set \(f(x)=1+x^2\):
  \[
    \int_0^1dx\,\frac{1+x^2}{(1-x)_+}
      =\int_0^1dx\,\frac{x^2-1}{1-x}
      =-\int_0^1(x+1)\,dx=-\frac32.
  \]
  For the first moment the test function is \(x(1+x^2)\):
  \[
    \int_0^1dx\,x\frac{1+x^2}{(1-x)_+}
      =-\int_0^1(x^2+x+2)\,dx=-\frac{17}{6}.
  \]
  A suitable \(\delta(1-x)\) term represents virtual corrections and
  adjusts these moments so that the required number or momentum
  conservation relation is exactly satisfied.

  \SupplementarySolution{22}{Mellin Moments Diagonalize a Convolution}
  Let \(M_n[f]=\int_0^1dx\,x^{n-1}f(x)\).  Then
  \[
  \begin{aligned}
    M_n[P\otimes f]
      &=\int_0^1dy\,\frac{f(y)}y
        \int_0^y dx\,x^{n-1}P(x/y)\\
      &=\left[\int_0^1dz\,z^{n-1}P(z)\right]
        \left[\int_0^1dy\,y^{n-1}f(y)\right].
  \end{aligned}
  \]
  Thus
  \[
    \boxed{M_n[P\otimes f]=M_n[P]M_n[f].}
  \]
  The integro-differential DGLAP equation becomes, for each \(n\), an
  ordinary differential equation in the scale; in the singlet sector it
  is a finite-dimensional matrix equation.

  \SupplementarySolution{23}{Conservation Constraints on Splitting Kernels}
  Nonsinglet flavor number is the zeroth moment of a valence
  distribution.  Its evolution vanishes only if
  \[
    \boxed{\int_0^1dz\,P_{\rm NS}(z)=0.}
  \]
  Total momentum is the first \(x\)-weighted singlet moment.  For each
  parent parton \(j\), conservation requires
  \[
    \boxed{\sum_i\int_0^1dz\,zP_{ij}(z)=0.}
  \]
  Real-emission kernels are positive over \(0<z<1\) and generally have
  endpoint divergences.  Plus prescriptions and negative virtual terms
  proportional to \(\delta(1-z)\) supply precisely the contributions
  needed to make these conservation moments vanish.

  \SupplementarySolution{24}{Recognizing Higher Twist}
  The leading term has twist two.  Since a twist-\(\tau\) contribution is
  suppressed by \(Q^{-(\tau-2)}\),
  \[
    \boxed{\frac{H}{Q^2}\ \text{is twist four},\qquad
      \frac{K}{Q^4}\ \text{is twist six}.}
  \]
  They matter at modest \(Q^2\) because the power denominators are then
  not small.  Near \(x=1\), the invariant mass of the unobserved hadronic
  state decreases and the leading-twist phase space is squeezed;
  multiparton correlations and target-mass effects are consequently
  enhanced relative to the leading contribution.

  \SupplementarySolution{25}{A Factorial Series and Its Borel Pole}
  Direct summation gives
  \[
    B(z)=K\sum_{n=0}^\infty\left(\frac za\right)^n
      =\boxed{\frac{K}{1-z/a}},
  \]
  with a pole at \(z=a\).  If \(a>0\), the pole lies on the integration
  contour \(z\in[0,\infty)\), so upper and lower contour prescriptions
  differ.  If \(a<0\), the singularity is on the negative axis and the
  positive-coupling integral is unambiguous, although the perturbative
  series still has zero radius of convergence.

  \SupplementarySolution{26}{Alternating and Fixed-Sign Divergence}
  For \(f_n\sim(-1)^nn!A^n\),
  \[
    B(z)\sim\sum_n(-Az)^n=\frac1{1+Az},
  \]
  so the leading singularity is \(z=-1/A\).  For
  \(f_n\sim n!A^n\),
  \[
    B(z)\sim\frac1{1-Az},
  \]
  with a singularity at \(z=+1/A\).  The alternating series is Borel
  summable along the positive axis; the fixed-sign series requires a
  lateral or principal-value prescription because its pole lies on the
  integration contour.

  \SupplementarySolution{27}{An Ambiguity from a Positive-Axis Pole}
  The two lateral contours differ by a closed contour around the pole.
  Since the residue of \(B(z)=R/(z_0-z)\) at \(z_0\) is \(-R\),
  \[
    (gF)_+-(gF)_-
      =-2\pi iR\,e^{-z_0/g},
  \]
  up to the orientation convention.  Thus the magnitude of the
  ambiguity in \(F\) is
  \[
    \boxed{\left|\Delta F\right|
      =\frac{2\pi|R|}{g}e^{-z_0/g}.}
  \]
  No power series in \(g\) can reproduce the exponential, which is why
  perturbation theory alone cannot select the contour.

  \SupplementarySolution{28}{Renormalon Position and OPE Power}
  In the conventional QCD Borel variable, a singularity at \(u=p/2\)
  produces an exponential factor
  \[
    \exp\!\left[-\frac{u}{\beta_0\alpha(Q)}\right]
      =\exp\!\left[-\frac{p}{2\beta_0\alpha(Q)}\right]
      \sim\left(\frac{\Lambda}{Q}\right)^p.
  \]
  Therefore its contour ambiguity has precisely the size
  \[
    \boxed{\Delta C_{\rm pert}\sim(\Lambda/Q)^p.}
  \]
  Consistency requires an OPE operator whose matrix element supplies a
  power correction of the same order.  The Wilson-coefficient and
  matrix-element prescription dependences can then cancel.

  \SupplementarySolution{29}{Optimal Truncation of an Asymptotic Series}
  Consecutive terms have ratio
  \[
    \frac{T_{n+1}}{T_n}\simeq(n+1)A\alpha.
  \]
  Their magnitude decreases until
  \[
    \boxed{n_*\simeq\frac1{A\alpha}.}
  \]
  Stirling's formula gives a smallest term, up to a power prefactor, of
  order
  \[
    |T_{n_*}|\sim
      \exp\!\left[-\frac1{A\alpha}\right].
  \]
  The nearest Borel singularity is at \(z_0=1/A\), so this is
  \(e^{-z_0/\alpha}\), the same nonperturbative scale as the lateral
  Borel ambiguity.  Adding terms beyond \(n_*\) worsens the approximation.

  \SupplementarySolution{30}{Cancellation of a Renormalon Ambiguity}
  Write a physical quantity schematically as
  \[
    \mathcal A=C_0^{(\mathcal P)}(Q)
      +\frac{C_d(Q,\mu)}{Q^d}
       \langle O_d(\mu)\rangle_{\mathcal P}+\cdots,
  \]
  where \(\mathcal P\) labels a Borel/subtraction prescription.  Changing
  it shifts the perturbative coefficient by
  \(\delta C_0\sim(\Lambda/Q)^d\).  The separation between short and long
  distances also shifts the renormalized matrix element so that
  \[
    \boxed{\delta C_0+
      \frac{C_d}{Q^d}\delta\langle O_d\rangle=0.}
  \]
  Only the sum is physical.  A perturbative coefficient or condensate
  quoted separately is therefore incomplete unless its common
  prescription is specified.
\end{enumerate}
